\begin{textsample}{POS Dim 2 – vem\_ed\_20 – Score 4.00 – vem\_ed\_20/t101.txt}  \label{ex:f2_pos_101}
“Para a longevidade dos projetos, é importante ter a crença nos benefícios que trazem para professores, alunos e instituição; mostrar-se confiável ao parceiro por meio de uma comunicação intensa e honesta; ter flexibilidade, compreensão das características culturais do colega estrangeiro e otimismo”, resumiu Moreira.
Na última \textbf{década}, os PCIs envolveram 50 Fatecs, 7.484 alunos e 49 instituições de ensino superior em 17 países.
Entre 2018 e 2022, a Unesp realizou 95 projetos em 27 unidades com 54 instituições parceiras em 17 países, envolvendo 1.924 alunos. “Para garantir a longevidade do programa de Intercâmbios Virtuais, é importante informar a comunidade acadêmica, oferecer formação de docentes, auxiliar na \textbf{busca} de parceiros internacionais, reconhecer as ações de internacionalização no âmbito institucional e registrar as atividades de Intercâmbio Virtual nos históricos escolares dos alunos”. “Intercâmbios Virtuais e extensão com a comunidade” foram abordados pelos professores Jan Krimphove (coordenador de assuntos internacionais da Unichristus / Ceará) e Silvio Ribeiro (professor da Fatec Lins).
A universidade cearense desenvolve colaborações com a DePaul University (EUA) na área de saúde. “Se bem \textbf{planejados}, os Intercâmbios Virtuais podem \textbf{contribuir} para as atividades de extensão; além disso, são importantes oportunidades de internacionalização para desenvolver a cidadania global dos estudantes”, comentou Jan Krimphove.

% matched lemmas: busca, contribuir, década, planejado
\end{textsample}
